% Replacement for Section II, subsections B--D only.
% Retain the paper's existing Section II introduction, subsection A,
% and subsection E onward. No additional packages are required.

\subsection{Digital Text (``Papers'') Workflow}
\label{sec:paper-workflow}

We construct a schema from a publication using only evidence in that
publication. Dataset files, repository documentation and human annotations
are excluded from paper-side prompts. A separate workflow extracts a schema
from the released dataset (Section~II-C). Both use the
same schema representation and produce two outputs: a schema for comparison
and a separate evidence record for inspecting its sources and decisions.

The paper workflow establishes names and table membership before adding
attributes. We describe it in three stages. A \emph{stage} contains related
steps; a \emph{step} is one minimal operation, such as reading a passage,
requesting a decision or checking a response. Review can return a stage to
its source evidence for bounded repair, as described in
Section~II-D.

\subsubsection{Stage 1: Prepare passage evidence:}
We extract the PDF text and record the source identity, physical page
numbers and character intervals. These locators allow later claims to be
checked against the original input. A relevance index helps locate passages
that may describe a schema. Uncertain passages are treated conservatively
as relevant. The index guides navigation; name discovery still scans all
source passages, including those marked irrelevant.

\subsubsection{Stage 2: Discover names and group tables:}
The LLM proposes candidate dataset views, fields and descriptive elements
from the paper. It then keeps, merges or excludes candidates using the
entities represented by the records, their level of detail, and the files
and relationships described in the source. A printed table may describe part
of a dataset without being a
dataset table itself, and one dataset view may be described in several
passages. Grouping resolves these distinctions before attribute enrichment.
The result is an inventory of retained names and their table membership,
with supporting evidence.

\subsubsection{Stage 3: Enrich attributes:}
For small batches of established objects, the LLM receives focused source
passages and fills the requested attributes, such as datatype, unit or
encoding. Each attribute requires its own supporting evidence or an explicit
unknown decision. A supported name does not establish its datatype or unit.
Unsupported optional attributes remain unknown rather than being completed
from general knowledge or the released dataset. Because a focused request
may receive only selected passages, an unknown value does not prove that
the information is absent from the full paper.

\subsection{Schema Extraction from Dataset}
\label{sec:dataset-workflow}

The dataset workflow uses the released files and their supporting
documentation. Native adapters read structural facts, while an LLM
interprets their meaning and enriches the schema. Schema assembly is
deterministic. As on the paper side, names and structure are established
before attributes are filled, and the schema and its evidence are retained
separately.

\subsubsection{Stage 1: Read native facts and fix the sample:}
Adapters record the physical column inventory, column order and available
native type metadata. We preserve the original source bytes and retain a
fixed sample of at most 200 rows per selected table or file using random
seed 42, keeping all rows when fewer are available. Every subsequent step uses that same saved
sample. Claims based on sampled values therefore describe those observations
and do not establish constraints over the entire dataset.

\subsubsection{Stage 2: Establish names and bind documentation:}
Physical columns determine the field inventory and order. Position-based
identifiers distinguish columns with duplicate names. Supporting
documentation is associated with the corresponding tables and fields, so
semantic decisions can refer to both native facts and relevant explanations.
The resulting inventory fixes the objects available for attribute filling.

\subsubsection{Stage 3: Enrich attributes:}
The LLM uses the native readout, fixed sample and associated documentation
to interpret fields and fill the configured attributes. Evidence records
distinguish explicit declarations, direct observations and inferences.
For example, an observed string value is different from a declared type
constraint, and a unit inferred from a name is different from an explicitly
documented unit. Unsupported optional attributes remain unknown. The shared
harness reviews these decisions and assembles accepted outputs.

\subsection{Workflow Harness Details}
\label{sec:workflow-harness}

The harness coordinates source access, model requests, review and export
for both workflows. It preserves the evidence behind names, table membership
and attributes so that the final schema can be inspected at each level.

\subsubsection{Contract-bound requests and review:}
Model requests use a fixed JSON output contract specifying the required
decisions and evidence. Attribute responses must contain
the requested decision keys and refer only to objects in the established
inventory. Missing keys and newly introduced names are rejected. Accepted
claims retain a value, provenance status, reason and source references.
Provenance distinguishes \texttt{declared}, \texttt{observed} and
\texttt{inferred} information; unresolved decisions are exported as
\texttt{unknown}.

A separate review checks proposed decisions against their evidence and
returns feedback for bounded repair. Final checks cover names, structure and
claims; failed checks block acceptance. Unresolved optional attributes and
deferred source references remain visible. The harness retains source views
and model responses, allowing a reader to inspect both the selected values
and the decisions behind them. Contract and citation checks establish that
an output is well formed and source-linked; they do not guarantee that its
meaning is correct.

\subsubsection{Schema and evidence outputs:}
Accepted decisions are assembled deterministically into
\texttt{full-schema.json}. The schema describes datasets, tables and fields,
with structure, value and encoding attributes. This is the representation
used for subsequent paper--dataset comparison.
The separate \texttt{evidence.json} retains supporting evidence, conflicts,
coverage and unresolved decisions. Keeping both outputs makes the result
useful as a schema and as an inspectable record of how it was obtained.

\par\smallskip\noindent
\begin{minipage}{\columnwidth}
\setlength{\parindent}{1em}
The compact example below selects saved records from an accepted MIT
Supercloud dataset extraction. The outer grouping keys are for presentation;
\texttt{support\_excerpt} abbreviates the saved name quotation. Other shown
keys and values are retained from the records.

\begingroup
\footnotesize
\begin{verbatim}
{
  "schema_field": {
    "name": "Node",
    "value": {"datatype": "string"}
  },
  "retention": {
    "field_handle": "f0001",
    "table_handle": "table", "native_keep": true
  },
  "evidence": {"Node": {
    "name": {
      "status": "declared",
      "location": {
        "section": "field-0001.json#page-1",
        "physical_page": 1,
        "start_char": 0, "end_char": 1500
      },
      "support_excerpt": "\"name\":\"Node\""
    },
    "value.datatype": {
      "status": "observed",
      "location": {
        "section": "field-0001.json#row-1",
        "physical_page": 1,
        "start_char": 285, "end_char": 303
      },
      "support": "\"r2652301-n830961\""
    }
  }}
}
\end{verbatim}
\endgroup

The field's retained table membership, declared name and observed datatype
have separate records. Here page 1 is a page of the generated dataset
readout. The string quotation supports the sample-based datatype decision;
it does not establish a population-wide constraint.
\end{minipage}
